\scanpage{59}


\mspage{2}

La situation
\diagram{\begin{tikzcd}[ampersand replacement=\&,column sep=large,row sep=large]
\biggl\{\, C''\arrow[r,"j_*"] \& C\arrow[l,shift right=5,"j^*"']\arrow[l,shift left=5,"j^!"]\arrow[r,"\deleted{\unsure{i_*}}i^*"] \& C'\arrow[l,shift right=5,"i_!"']\arrow[l,shift left=5,"i_*"]\\
H \& F \& G
\end{tikzcd}}

(\uline{catégories abéliennes}, \uline{foncteurs additifs})

\[
\left\{\begin{aligned}
\operatorname{Hom}(F,j_*H)&=\operatorname{Hom}(j^*F,H)\\
\operatorname{Hom}(j_*H,F)&\simeq\operatorname{Hom}(H,j^!F)\\
\operatorname{Hom}(F,i_*G)&\simeq\operatorname{Hom}(i^*F,G)\\
\operatorname{Hom}(i_!\deleted{\unsure{F}}G,F)&\simeq\operatorname{Hom}(G,i^*F)
\end{aligned}\right.
\]

\[
\begin{aligned}
j_!=j_*,\,j^*,\,i^*=i^!,\,i_!&\quad\text{exacts}\\
j^!,\,i_*&\quad\text{exacts à gauche}
\end{aligned}
\]

\[
\varphi=j^*i_*:\ C''\longleftarrow C'\qquad\varphi\ \text{exact à gauche}
\]

\[
\left\{\begin{aligned}
j_*j^!F&=\underline{\Gamma}_Y F& j_*j^*F&=F^Y\\
i_*i^*F&=\underline{\Gamma}_U F& i_!i^*F&=F^U
\end{aligned}\right.
\]

\[
\left\{\begin{aligned}
i^*j_*&=0&\quad &\times\\
j^!i_*&=0&&\times\\
j^*i_!&=0&&\times\\
j^!i_!&=0
\end{aligned}\right.
\]

\[
j^*\longrightarrow\varphi i^*
\]

\uline{Prop. 1} Partons de $C''\xrightarrow{j_*}C\xrightarrow{i^*}C'$, où $j_*$ est pleinement fidèle, $i^*$ est un passage au quotient exact, et $j_*\simeq\operatorname{Ker}i^*$. \inserted{donc $i^*j_*=0$,} Alors

(i) Si $i_!$ existe\inserted{ [néc. pl. fid]}, alors $j^*$ existe\inserted{ [néc. un passage au quotient]}, et on a la suite exacte

\diagram{\begin{tikzcd}[ampersand replacement=\&,column sep=small]
C''\arrow[r,shift right,"j_*"'] \& C'\arrow[l,shift right,"j^*"']\arrow[r,shift right,"i^*"'] \& C''\arrow[l,shift right,"i_!"']
\end{tikzcd}}

\[
\boxed{i_!i^*F\longrightarrow F\longrightarrow j_*j^*F\longrightarrow 0}
\]

\[
\boxed{\deleted{\unsure{j^*i_!=0}}}
\]

De plus
\[
\left\{\begin{aligned}
j^*j_*&\overset{\sim}{\longleftrightarrow}\mathrm{id}_{C''}\\
i^*i_!&\xrightarrow{\sim}\mathrm{id}_{C'}\\
j^*i_!&\simeq 0_{C''}
\end{aligned}\right.
\]

(i bis) Si $i_*$ existe, alors $j^!$ existe, et on a
\[
i_*i^*F\longleftarrow F\longleftarrow j_*j^!F\longleftarrow 0
\]

De plus
\[
\left\{\begin{aligned}
j^!j_*&\xleftarrow{\sim}\mathrm{id}_{C''}\\
i^*i_*&\xrightarrow{\sim}\mathrm{id}_{C'}\\
j^!i_*&\simeq 0_{C'}
\end{aligned}\right.
\]

\margin{\uline{Question} Si $j^*$ existe est-il vrai que $i_!$ existe ?}

(N.B. En général, on ne peut rajouter des zéros\ldots)

% archive sheet number: 58
